% ---------------------------------------
%
%    Beispieldiplomarbeit
%
%    - text kann gel{\"o}scht, und mit eigenen Inhalten gef{\"u}llt werden. 
%	In dieser Datei kann die gesamte Konfiguration durchgeführt werden.*
%	
%
%    *mit Ausnahme der spezifischen Optionen in natbib.cfg, die nicht von Interesse sein sollten
% ---------------------------------------

\documentclass[
    smallheadings,  % kleinere {\"U}berschriften
    twoside,        % einseitig, nur rechte seiten
    liststotoc,     % listen in inhaltsverzeichnis aufnehmen
    bibtotoc,       % literaturverzeichnis in inhltsvz. aufnehmen
    headsepline,     % trennlinie unter kopfzeile
    12pt	%Schriftgröße
    ]{scrbook}

\usepackage{a4}  %a4 Seitenformat benutzen
\usepackage[left=4cm,right=2.5cm,top=3.7cm,bottom=4cm]{geometry}
\usepackage[english, ngerman]{babel} %Verwende deutsche und amerikanische Silbentrennung
\usepackage[utf8]{inputenc} %damit k{\"o}nnen Umlaute ganz normal geschrieben werden. 

\usepackage{listings} %Fuer Codelistings

\usepackage{subfigure} %f{\"u}r mehrteilige Grafiken
\usepackage{epsfig}    %damit funktioniert das einbinden von grafiken {\"u}ber epsfig.
\usepackage{graphicx}     % zum einbinden von grafiken
\graphicspath{{grafiken}{../}{kapitel}} %da sind m{\"o}gliche bilder fuer den includegraphics-Befehl zu finden (man muss dann nicht den ganzen Pfad bei includegraphics angeben. 

\usepackage{multirow}     %fuer kompliziertere Tabellen
\usepackage{longtable}	%fuer kompliziertere Tabellen
\usepackage{rotating}	%enable landscape pages
\usepackage{framed}
\usepackage{scrpage2}     % paket f{\"u}r kopf- und fu{\ss}zeilen
\pagestyle{scrheadings}   % kopzeilenseitenstil

%\usepackage{ifpdf} %provides the switch ``ifpdf'' to determine if latex or pdflatex is executed
%see ftp://ftp.ctan.org/tex-archive/macros/latex/contrib/oberdiek/ifpdf.pdf or doc folder

\usepackage{float} % adds location parameter [H] to [htb] to force figures and tables to a location, no floating

%%%% use hyperlinks and adjust settings, see http://en.wikibooks.org/wiki/LaTeX/Hyperlinks and
% ftp://tug.ctan.org/tex-archive/macros/latex/contrib/hyperref/doc/manual.pdf
\usepackage[
a4paper=true, % a4
pdftitle={The title of this paper}, %title to display for pdfs
pdfauthor={The Author}, %set author
%dvipdfmx, % must be used if pdf is created from the dvi file rather than by running pdflatex. Otherwise links won't work
colorlinks,% set link color for all links to black, i.e. invisible but working links
citecolor=black,%
filecolor=black,%
linkcolor=black,%
urlcolor=black, %
plainpages=false, %used with pdfpagelabels to make pdf readers show roman and arabic page numbers correctly
pdfpagelabels, %used with plainpages=false to make pdf readers show roman and arabic page numbers correctly
%pagebackref=true %false: no back-references in bibliography to citations in text
]{hyperref} %include hyperlinks for citations, urls, and references
%CAUTION: breaks color settings in dvi, pdf works fine
%%%%end: use hyperlinks

\usepackage{setspace}
\usepackage{url}          % fuer urls: schreibweise ist z.B. \url{http://www.uni-mannheim.de}
%\urlstyle{same} % makes urls in the same font as the rest of the text


%Automatisch Abkürzungsliste generieren: 2x latex, makeindex myfile.nlo -s nomencl.ist -o myfile.nls, 2x latex
\usepackage{nomencl} % um Abkürzungen aufzunehmen: \abbrev{PDA}{personal digital assistant}
\let\abbrev\nomenclature
\renewcommand{\nomname}{List of Abbreviations} %Titel der Liste hier ändern, eine von zwei Optionen nutzen
%\renewcommand{\nomname}{Abkürzungsverzeichnis} %Titel der Liste hier ändern, eine von zwei Optionen nutzen
\setlength{\nomlabelwidth}{.25\hsize}
\renewcommand{\nomlabel}[1]{#1 \dotfill}
\setlength{\nomitemsep}{-\parsep}
\makenomenclature
\newcommand{\Listofabbrev}{
\printnomenclature
\newpage
}



%Inhaltsverzeichnis
\usepackage
%[
%tocfullflat,			%alle Eintraege left-aligned, alternativ: tocflat
%tocbreakscareless		%page break between toc entry allowed
%]
{tocstyle}			%mehr Kontrolle ueber Inhaltsverzeichnis
\usetocstyle{allwithdot}	%ensure there are dots in table of contents, alternativ: noonewithdot, nopagecolumn
% docu http://www.tug.org/texlive/devsrc/Master/texmf-dist/source/latex/koma-script/tocstyle.dtx
\setcounter{secnumdepth}{3} %numbering up to fifth level in table of contents
\setcounter{tocdepth}{3} %show entries in toc up to fifth level (subsubsubsubsection)


\setlength{\parindent}{0pt}		% Einzug fuer neuen Absatz
\setlength{\parskip}\medskipamount % Abstand neuer Absatz, besser als explizite Angabe in pt

\onehalfspacing %Zeilenabstand 1,5


% kapitel{\"u}berschriften in schriftart mit serifen
\setkomafont{sectioning}{\normalfont\normalcolor\bfseries}

% gestaltung der kopfzeilen
\ohead{\pagemark}
\ifoot{}
\cfoot{}
\ofoot{} 
\cohead{}
\ihead{\headmark}
\setkomafont{pagehead}{\normalfont\bfseries}
\setkomafont{pagenumber}{\normalfont\bfseries}
\automark{section}

% ----- ende der pr{\"a}ambel ----------------------------------






\begin{document}  % dokument f{\"a}ngt an
\selectlanguage{english} %englische Silbentrennung, fuer deutsche Arbeiten: \selectlanguage{ngerman}
\frontmatter      % vorspann, kapitel r{\"o}misch nummeriert

\newgeometry{margin=3cm}
% Die Titelseite der Arbeit

\begin{titlepage}

\begin{center} % zentrieren

  % Logo der Universit{\"a}t Mannheim
  \begin{figure}[ht]
    \centering
    \includegraphics[width=.6\textwidth]{grafiken/unilogo.png}
  \end{figure}
  
  % Vertikaler Zwischenraum
  \bigskip
  \vfill 
  %\begin{framed}
  % Titel der Arbeit und Typ der Arbeit, umrandet
    \begin{center}
     \textsc{{\LARGE Enhancing Sustainable Software Engineering Approaches / Models through SECoMo\\}}
                                % Letztes \\ ist wichtig, beginnt eine neue Zeile f{\"u}r die Art der Arbeit
  
      \bigskip
  
                                % Art der Arbeit, ggf. auszutauschen gegen Seminar- oder Doktorarbeit
      \textbf{Seminar Thesis}
    \end{center}
   % \end{framed}
    \vfill
    \vfill
  
  % Daten des Erstellers, Einreichungsdatum
  % in einer Tabelle ausgerichtet
  \begin{tabular*}{0.62\textwidth}{r@{\extracolsep{\fill}}l}
   submitted: &\ May 2017\\\\
    by: &\ Natalie Buchner\\
		&\ nbuchner@mail.uni-mannheim.de\\
    &\ born July 12th 1993\\
    &\ in Bonn\\
    \\
    Student ID Number: &\ 1496726\\
  \end{tabular*}
  \vfill
  \vfill
  
  % Unten: Kontaktdaten des Lehrstuhls f{\"u}r Softwaretechnik
  
  \rule{\textwidth}{.4pt}\\ % vertikale Linie
  University of Mannheim\\
  Chair of Software Engineering\\
  D -- 68159 Mannheim\\
  Phone: +49 621-181-3912, Fax +49 621-181-3909\\
  Internet: \url{http://swt.informatik.uni-mannheim.de}
\end{center}

\end{titlepage} % Ende des Titelblatts

%%% Local Variables: 
%%% mode: latex
%%% TeX-master: "~/Documents/DA-Vorlage/beispiel/da-beispiel"
%%% End: 
     % titelseite einbinden
\restoregeometry
\chapter{Abstract}

%TODO - write abstract & send to Atkinson  % abstract einbinden
\tableofcontents            % inhaltsverzeichnis

\listoffigures              % abbildungsverzeichnis
\listoftables               % tabellenverzeichnis

\clearpage %tell toc that there is a pagebreak after list of tables, however the hyperlink still points to the wrong page
\addcontentsline{toc}{chapter}{\nomname} %abkuerzungen ins inhaltsverzeichnis

\Listofabbrev % liste der abkuerzungen erstellen



\mainmatter       % hauptteil, kapitel lateinisch nummeriert
%\include{kapitel/content}
\chapter{Introduction}
%Introduction into the topic
%s. Abstract, nur ausführlicher?!

%description of work done (more detailed)
%(... später, wenn auch tatsächlich was geschrieben wurde...)

%Structure of the thesis
%(...)





%\chapter{Schriftliche Ausarbeitung}
%
%\section{Formale Anforderungen}
%
%Der nachstehende Abschnitt gibt einen kurzen Überblick über die
%formalen Anforderungen an die schriftliche Ausarbeitung von Seminar-,
%Studien-, Bachelor-Abschluss- und Diplomarbeiten. Im Weiteren bitten wir
%unbedingt, die für Sie gültige Prüfungsordnung zu beachten.
%
%Für alle Ausarbeitungen sind die Templates des Lehrstuhls (\url{http://swt.informatik.uni-mannheim.de/de/studies/master-and-bachelor-theses/guidelines/}) zu verwenden.
%
%\begin{table}[tb]
%\centering
%\begin{tabular}{|p{0.35\textwidth}*{1}{|p{0.25\textwidth}}|p{0.25\textwidth}|}
%\hline
%& \textbf{Bachelorarbeit} & \textbf{Masterarbeit}\\
%\hline
%Seitenzahl (ohne Anhang) & ca. 30 & ca. 80\\ 
%\hline
%Abbildungsverzeichnis & ja & ja\\
%\hline
%Tabellenverzeichnis & ja & ja\\
%\hline
%Ehrenwörtliche Erklärung & ja & ja\\
%\hline
%Abtretungserkl{\"a}rung & nach Vereinbarung & nach Vereinbarung\\
%\hline
%Bearbeitungszeit & max. 3 Monate & max. 6 Monate\\
%\hline
%Anzahl der abzugebenden Examplare (+ CD) & 3 Exemplare (gebunden) an LS Sekretariat & 3 Exemplare (gebunden) an LS Sekretariat\\
%\hline
%\end{tabular}
%\caption{Übersicht der Anforderungen \label{tableAnforderungen}}
%\end{table}
%
%Bei Abweichungen von diesen formalen Anforderungen, insbesondere bei der Seitenanzahl, ist unbedingt Rücksprache mit dem zuständigen Betreuer zu halten. Bei der Anfertigung der Arbeit sind die Grundregeln des wissenschaftlichen Arbeitens zu beachten. Zum Beispiel sind Quellen anzugeben, wörtlich aus Quellen übernommene Textstellen als Zitate zu kennzeichnen usw.
%
%\subsection{Verwendete Literatur}
%
%Gute Literatur zur Bearbeitung der Arbeit kann unter folgenden Adressen gefunden werden:
%
%\begin{itemize}
%\item IEEE Computer Society Digital Library: \url{http://www.computer.org/portal/site/csdl/index.jsp}
%\item ACM Digital Library: \url{http://portal.acm.org/dl.cfm}
%\item CEUR Workshop Proceedings: \url{http://ceur-ws.org/}
%\item Springer Link: \url{http://www.springerlink.com/home/main.mpx}
%\end{itemize}
%
%Die Dokumente von ACM als auch Springer sind mit einer universitären IP (z.B. über VPN-Verbindung oder aus dem Pi-Pool) kostenlos zu bekommen. Zu IEEE besitzt die Universitätsbibliothek Zugangsdaten. CEUR dagegen enthält nur frei zugängliche Dokumente.
%
%Das Literaturverzeichnis ist enstprechend den Vorgaben des Templates aufzubauen. Ein Beispiel f\"ur eine Zitation: \cite{tannenbaum}.
%
%\subsection{Inhalt}
%
%Interessant Hinweise zum Erstellen einer Forschungsarbeit liefert z.B. \url{http://www.cs.cmu.edu/~Compose/shaw-icse03.pdf}
%
%\subsection{Abgabe}
%
%Die schriftlichen Ausfertigungen müssen fristgerecht im Sekretariat oder beim Betreuer abgegeben werden. Ansonsten kann nach den Bestimmungen der Prüfungsordnung keine erfolgreiche Leistung bescheinigt werden. 
%
%Jede Arbeit soll doppelseitig gedruckt werden und dem Layout eines guten
%Buches entsprechen. Wie in der Tabelle oben ersichtlich, benötigen werden drei gebundene Exemplare. Alle Exemplare der Bachelor-Abschluss- bzw. Diplomarbeit müssen die ehrenwörtliche Erklärung unterschrieben beinhalten. Auf der inneren Rückseite von jedem Exemplar ist eine eingebundene/eingeklebte Hülle mit einer CD mit abzugeben. Der Inhalt dieser CD:
%
%\begin{itemize}
%\item letzte Fassung der Arbeit (auch in veränderlicher Form, z.B. *.doc, *.tex)
%\item Foliensatz des Abschlussvortrags (auch in veränderlicher Form, z.B. *.ppt, *.tex)
%\item Quellcode der programmierten Anwendung
%\item ausführbares Programm (und eine Anleitung, wie dieses zu verwenden ist)
%\item Dokumentation zu dem Programm (z.B. Javadoc)
%\item Referenzen (in der Arbeit verwendete Dokumente, die in elektronischer Form vorliegen)
%\end{itemize}
%
%Der Schein f{ür die Seminararbeit bzw. der erfolgreiche Abschluss der Studien-, Bachelor-Abschluss- bzw. Diplomarbeit wird dem Studienbüro erst gemeldet, wenn die Arbeit ordnungsgemäß abgegeben und der Vortrag absolviert wurde. 
%
%Eventuell ausgehändigte Schlüssel, Bücher oder sonstige Leihgaben des Lehrstuhls müssen bei der Abgabe der schriftlichen Ausfertigung beim Entleiher bzw. im Sekretariat zurückgegeben werden.
%
%\subsection{Abschlussvortrag}
%
%Kandidat/in und Betreuer vereinbaren einen Termin für den Vortrag, der in Anwesenheit von Herrn Prof. Atkinson gehalten wird. Die Folien sind dem Betreuer eine Woche vorher zuzusenden. Eine wichtige Fähigkeit bei Präsentationen (z.B. bei Konferenzen, berufliche Präsentationen) ist, Inhalte in begrenzter Zeit auf den Punkt zu bringen. Daher sollte die Vortragsdauer unbedingt eingehalten werden (Nichteinhalten kann zu einer schlechteren Benotung führen).
%
%\begin{table}[h]
%\centering
%\begin{tabular}{|l|c|c|}
%\hline
%& \textbf{Bachelor} & {\textbf{Masterarbeit}}\\
%\hline
%Vortragsdauer & 15min & 20min\\
%\hline
%Diskussion & 5min & 10min\\
%\hline
%\end{tabular}
%\caption{Vortragsdauer \label{tableVortragsdauer}}
%\end{table}
%
%%In this section, there is a first picture \ref{flow1}.
%
%\begin{figure}[h]
%\centering
%\includegraphics[width=0.07\textwidth]{grafiken/grafik}
%\caption{Beispielgraphik}
%\label{flow1}
%\end{figure}

\chapter{Background} %or Foundations?!
%What is used in the thesis e.g. formalisms, methodologies, ... (@ atkinson - any formalism / methodologies needed?!)
%Clearly separate between contributions (by me) and foundations (from others)
%Only describe what is needed to understand the thesis

\section{Sustainability and Sustainable Software Engineering}
% Definitions - Sustainability, Sustainable Software & Sustainable SW Engineering

% Impacts of software - why is sustainable sw engineering important?

\section{Related work in Sustainable Software Engineering}
% RELATED WORK - What exists?
\subsection{Principles and Manifestos}
% principles, design, guidelines etc. (...??????) karlskrona manifesto!?
% in which aspects do they occur - strategic level for company, in design of SW & architecture, 
% considered in requirements engineering - what needs to be fulfilled? what aspects occur? (see contributions of papers / etc.)

\subsection{Models and metrics}
% mainly: what kind of general approaches exist / models?
% Greensoft - lifecycle model, impacts (but not - metrics!?)
% ...

% - or rather in the contribution part? critique - what is missing? - not used in practice?! (quellen?!) - why?

\subsection{The SECoMo approach}
% RELATED WORK / FOUNDATION!?Chapter on Secomo
%evt. Kobra?
% Green Specifications - Background Idea

% what does secomo do? what's used from green specifications?
% --> simple outline?!
\section{The SECoMo approach} % RELATED WORK / FOUNDATION!?
%�berleitung: vor allem: warum nehme ich das so raus? oder reicht es das in der einleitung zu erkl�ren?!?!
%�berleitung: very general - where does it fit in? - or NOT and just put it together in the "contributions"' part
A relatively new addition to the field of sustainable software engineering approaches %ok to say it like that?
 is the %wie Sachen hervorheben??
\textbf{Software Eco-Cost Model} %wie ging das nochmal - abk�rzungen einf�hren?
(SECoMo) Approach by Thomas Schulze %wie geht das bei Latex - nur Jahr erw�hnen oder so? - oder welchen zitierstil soll ich nutzen?
\cite{schulze_cost_2016}. This approach provides % wen? software engineers? theoretisch auch andere!
Software Engineers with generic models and metrics necessary to estimate and express the ecological costs %(of the use of) 
a software system causes when it is used \cite{schulze_cost_2016}. Thus, SECoMo represents a concrete estimation approach for the impact a software system has regarding ecological sustainability during its usage phase. 

\subsection{What is SECoMo about?} %% BILD einf�gen???
%vllt 1, 2 direkte zitatie mit einbinden?
The main motivation behind SECoMo is to provide an approach that allows to not only measure the ecological costs that are actually caused by a software system, but also to be able to estimate those costs upfront, for example already during the design phase of a software engineering project \cite{schulze_cost_2016}. %%mention here: WHY!!! -> because it is cheaper to make changes here, thus it is better to have more information already in this pahsee!!!
 In order to achieve this, SECoMo offers a set of mathematical models which allow to precisely calculate eco-cost metrics, based on information that is already available in the design phase: specification models %????? %name the auxiliary models that are created from that already?  %what about the "`generic"' part?
that describe the functionality, behavior and structure of the software system \cite{schulze_cost_2016}. Furthermore, SECoMo is intended to be highly adaptable %??? other word,
in order to allow the estimates to be calculated for different levels of details available - from an early level where only very general information about the software system is available, over an intermediate level with partly more detailed information, to an advanced level with very specified details that allow for more accurate estimates \cite{schulze_cost_2016}.\\ %is that okay...? %
In addition to the mathematical models, the SECoMo approach also defines a set of eco-cost drivers in order to identify causes for certain ecological impacts a software system has and to better describe under which circumstances they occur \cite{schulze_cost_2016}. The auxiliary models used in SECoMo which extend the specification models %?????
provide information about these cost drivers, but can also be used to express the estimated eco-costs of the software system \cite{schulze_cost_2016}. This way, SECoMo additionally offers a possibility to communicate estimated or measured eco-costs to stakeholders of a project which can use this information to make improved decisions \cite{schulze_cost_2016}.\\ %what about: comprehensible and clear - on the right level!?
%TODO: %mention: how are Eco-Costs defined!? to give context %mention callibration process???
%what benefits does it offer -> pro's / con's, conclusion? - when to use it? - or put that at the end of the chapter? maybe not to much.. or this could be starting point for comparison - rewrite?!
Against this background, the SECoMo approach is intended to be used in the early stages of software engineering projects to create estimates about the ecological impact of a software system, so as to enable transparency about the sustainability aspect right from the start \cite{schulze_cost_2016}. This again makes it possible for software engineers and other stakeholders to make decisions about changes to the software at the design stage which take the impact on ecological costs into account %\cite{schulze_cost_2016}
- be it to improve certain eco-cost critical aspects of the software because ecological sustainability is a major concern, or to at least be aware of the eco-cost trade-offs other decisions cause that might be motivated by other concerns, e.g. profitability. \cite{schulze_cost_2016}\\ %???? possible?
%+++++++ the aspect that it is CHEAPER!!! to make changes in this phase!!!!!
But SECoMo can also be used in the context of defining requirements for a software system, for example in terms of specifying upper bounds for the eco-cost metrics that must not be exceeded, or even to calculate exact eco-costs if enough details are given, e.g. for the specific usage scenario of a software system in a certain environment %when all hardware details are known?!
\cite{schulze_cost_2016}.\\

\subsection{How can SECoMo be integrated in a software engineering project?}
%briefly: "`algorithmic"' estimation model? what does it consist of -metrics, (auxiliary models? -> how? not very specific?), & cost-drivers - other order?
%HOW does it work? how to "use" SECoMo? to have a "`process"' - see evaluation chapter
On a very high level, the SECoMo approach, when being used within a software engineering project in order to estimate the eco-costs, can be structured as follows:\\
In order to start with the SECoMo approach, the initial design phase of the software needs to be in a stage where it is already clear on some level how the software will be structured, which behavior it should follow and what functionalities it should provide. %What to you need to begin with? software needs to be designed at least to a certain level where information about the structure - what components / classes / data types? - and the behaviour & functionality of these components is available / already specified
 Ideally, this is reached by using a modeling approach that creates a structural, behavioral and functional view on the modeled software. SECoMo does not require a specific modeling approach that needs to be used to create these views, but these model types are the basis for the further steps of using SECoMo as estimation technique \cite{schulze_cost_2016}.\\
If this basis is given, the following steps of the SECoMo approach can be integrated already in the earliest stages of a software development process (if early eco-cost estimations are required), otherwise also in later iterations of the process to refine the initial estimates or to calculate the actual eco-cost values (for example in the testing phase):
\begin{enumerate}
	\item Calibration of the models, in order to receive accurate values concerning the hardware energy consumption factors (resource factors) that are later needed as inputs to the mathematical models \cite{schulze_cost_2016}; %and resource factors?! -> to implicitly cover these implementation % related things.. somehow
	%so that the information about the hardware aspects is accurate and relates to a given situation %???
	\item Preparation of the auxiliary models so that they include all information available regarding the relevant eco-cost drivers (on the intermediate and advanced level certain information can be derived using Markov Models) \cite{schulze_cost_2016}; %e.g. information needed about frequency - of operation, or how often data type is used in memory / persistent?!
	\item And finally Calculation of the estimates for the different eco-cost metrics based on the available information %-> so, which level of detail?
	and the specified parameter values, which identify aspects like user type or case \cite{schulze_cost_2016}.
\end{enumerate}
Based on the estimated or calculated eco-costs it is also possible to derive specific reports on the environmental impact of the software for different types of stakeholders, or to simply communicate the eco-costs via the existing auxiliary models \cite{schulze_cost_2016}. This way, it is possible for stakeholders to make informed decisions on the software design taking the eco-cost information into account, and especially doing so at an early level \cite{schulze_cost_2016}.
% e.g. make decisions, like - oh this data type causes a lot of eco-costs - could be the size of ceratin attributes
% or this scenario is highly (consumative?!) -> we should just decide not to allow certain combinations of actions to be possible


%%%%%%%%%%%%%%%%%%%%%%%%%%%%%%%%%%%%%%%%%%%%%%%%%%%%%%%%%%%%%%%%%%%%%%%%%%%%%%%%%%%%%%%%%%%%%%%%%%%%%%%%%%%%%%
%improve wording, e.g. - not "provide" and "support", "help", but better synonyms!
%Multiple sections, based on what I do
\chapter{Contributions}
\section{Why sustainable software engineering is not yet used in practice(?!)}
% of existing models / methods / metrics:
% critique - what is missing? - not used in practice?! (quellen?!) - why?

\section{SeCoMo in the context of Sustainable SW Engineering}
%What is done in the thesis?

% How do they relate? What does it contribute to sustainability?


% Find / create some ideas - where to integrate it in some models
% - in lifecycle models

% with metrics / AS metrics

% for eco-label?!


% !!! Integrate with GREEN-IT readiness / maturity models!?


\chapter{Findings(?)} %or above part here... - AND: separate from Conclusion?!
% What is enhanced? Answer the question here...


%\chapter{Limitations}
% of secomo - what it still does not contribute - social sust. (more a process related aspect?)
%+ in general: social sust. -> still not much work done in general (see paper Hinai & CHit..)

%of the integration - how it can't be perfectly integrated?!


%The main contributions / findings of the work
\section{Conclusion}
In the context of Sustainability, ICT plays a major role as contributor to negative impacts, but also as enabler to improve Sustainability - in the ICT sector itself, but also in other sectors. Even though it might not be obvious at the first glance, software products are relevant in that context, too: They can have direct and indirect impacts on the different dimensions of Sustainability, for example by influencing a hardware product's resource consumption or by promoting sustainable life styles.\\
The research field of sustainable software engineering deals with concepts and approaches on how to ensure that software products can improve their positive impact on Sustainability during their whole life cycle. Over the past few years, the amount of contributions to the field has significantly increased, and a variety of concepts exist that tackle different aspects of software sustainability. But a common understanding and concrete guidance is still missing in the field, and despite the variety of approaches, not much is evaluated in practice, yet.

The SECoMo approach introduced in this paper is yet another addition to the contributions of this field, in particular in the area of approaches dealing with software usage sustainability. %something more about its classification?
It has the ability to enhance existing approaches in the field in various ways and might thus be suitable to help leading Sustainable Software Engineering to more adaption in real life projects. In any way, SECoMo can play an important role in the field in different areas:\\
With regards to existing metrics, SECoMo provides a new set of eco-cost metrics that can be applied to various types of software development projects as they are more flexible and can cover a broad range of contexts.\\
In comparison to existing sustainability models in the field, an important role that SECoMo can play is to enhance the already well-rounded GREENSOFT model with a concrete estimation approach for environmental sustainability aspects in order to make its development procedure model more likely to be adapted in practice. Over all, SECoMo introduces a new type of approach to Sustainable Software Engineering with its ability to estimate eco-costs even before development phase has started.\\
SECoMo can even contribute to the important area of Requirements Engineering by providing a much needed vocabulary to define sustainability requirements and also the respective method how to assess them. %%%%Abstract: As SECoMo can be integrated in all development phases, especially the early ones, it can help to enhance existing life cycle models with a concrete method for understanding and improving ecological sustainability in the design and implementation phases of software engineering. In addition, with its new set of sustainability metrics, SECoMo offers new options for sustainability measurement and assessment in existing models and tools, which base on a general way of software specification that can reasonably be applied in practice.
%can it help in practice?
%limitations?

The future work on SECoMo will show whether it is indeed a suitable approach that can enable the adoption of Sustainable Software Engineering in practice, and in any way there is still some work to do in the research field itself, starting from finding a common understanding of the underlying concepts over to evaluating the existing variety of approaches in real life projects. But the trend of Green IT and Sustainability adaption in companies for example shows, that the work in this direction is important and very much needed.
%ausblick - secomo - needs to be tested in practice -> and in combination! to say more about practice - but good starting point as very good for all kinds of sw projects but also very adaptive for different needs of stakeholders
%LATER also: more work for other dimensions of sutainbility?! - even though according to penzenstadler for example, every other of her 5 dimensions is more covered than environmentall? whatever
%\selectlanguage{ngerman} % jetzt sprechen wir deutsch.
%\chapter{Implementierung}

Zusammen mit dem Betreuer werden use-cases entwickelt anhand deren die
Software programmiert werden soll. Diese dienen auch als Bewertungsgrundlage.\\
\\
Die allgemein empfohlene Verzeichnisstruktur eines Projektes sieht wie folgt
aus:

\begin{itemize}
\item projektname
  \begin{itemize}
  \item bin
  \item doc
  \item lib
  \item src
  \end{itemize}
\end{itemize}

Die zu erstellende Software soll im package
de.uni\_mannheim.informatik.swt.\texttt{projektname} unter \texttt{src} liegen.\\
\\
Bei der Programmierung sollte durchgängig die englische Sprache verwendet
werden. Hierzu zählen insbesondere Kommentare im Quellcode, Namen von
Funktionen, Variable, Menüpunkte im Benutzerinterface, kurze
Hilfestellungen und Ausgaben von Programmen.\\
\\
Der Code sollte mit Hilfe von \texttt{lstinputlisting} formatiert und ausgegeben werden, wie in folgendem Beispiel:

\lstinputlisting[
  language=Java, numbers=left, stepnumber=5, firstnumber=1, breaklines=true, 
  basicstyle=\footnotesize,
  numberstyle=\tiny,
  caption={GuestbookForm.java},
  captionpos=b,
  label=GuestbookForm
]
{code/GuestbookForm.java.txt}

%\selectlanguage{english} % jetzt sprechen wir wieder englisch
%\selectlanguage{ngerman} % wenn im Inhaltsverzeichnis Appendix stehen soll, muss Engl gewaehlt sein, fuer Anhang Deutsch
\appendix

%Bibliographie
%waehle einen der folgenden 4 Eintraege
%\bibliographystyle{literatur/natdin} %DIN Style Literaturverzeichnis, comment out pagebackref
%\bibliographystyle{literatur/IEEEtran} % IEEE Style Literaturverzeichnis
%\bibliographystyle{literatur/natdinCustomized} %DIN Style Literaturverzeichnis + Punkt hinter jeder Literaturangabe -> low level config fuer Zitate: natdin.cfg im Projektordner, comment out pagebackref
%\bibliographystyle{literatur/natdinCustomizedEnglish} %DIN Style auf Englisch getrimmt mit Punkt hinter Literaturangabe -> low lovel config fuer Zitate: natdin.cfg im Projektordner, pagebackref kann damit nicht genutzt werden
%\bibliographystyle{aer} % alternativ auch apalike, aer, apalike2... s.a. http://web.reed.edu/cis/Help/LaTeX/bibtexstyles.html
\bibliographystyle{plain} % very nice bib style
%note on aer: does not like inbook entries
%\bibliographystyle{natdin}

\bibliography{literatur/lit} %Pfad zur bib-Datei

%appendices can be defined here, the appendix structure has to be added manually to the toc (table of contents)
%\clearpage  %toc new page
%\addcontentsline{toc}{chapter}{Appendix} %add chapter to toc
\addpart{\appendixname}
\chapter{First class of appendices}
\section{Some appendix}

This is a sample appendix entry.

% \newpage
% \input{kapitel/methodologyQuestionMappingTable}
% 
% \input{kapitel/question_origins}
% \include{extern/fragebogenE}
% \include{extern/fragebogenD}
% 
% \input{kapitel/invitationLetters}
% 
% 
% \include{kapitel/coverageTable}

%Eidesstattliche Erklaerung
\chapter*{Eidesstattliche Erkl\"{a}rung}
\thispagestyle{empty}
Hiermit versichere ich, dass diese Abschlussarbeit von mir persönlich verfasst
ist und dass ich keinerlei fremde Hilfe in Anspruch genommen habe. Ebenso
versichere ich, dass diese Arbeit oder Teile daraus weder von mir selbst noch
von anderen als Leistungsnachweise andernorts eingereicht wurden. Wörtliche oder
sinn\-gemäße Übernahmen aus anderen Schriften und Veröffentlichungen in gedruckter
oder elektronischer Form sind gekennzeichnet. Sämtliche Sekundärliteratur und
sonstige Quellen sind nachgewiesen und in der Bibliographie aufgeführt. Das
Glei\-che gilt für graphische Darstellungen und Bilder sowie für alle
Internet-Quellen.

Ich bin ferner damit einverstanden, dass meine Arbeit zum Zwecke eines
Plagiatsabgleichs in elektronischer Form anonymisiert versendet und gespeichert
werden kann. Mir ist bekannt, dass von der Korrektur der Arbeit abgesehen werden
kann, wenn die Erklärung nicht erteilt wird.
\bigskip

\vspace{2.5cm}

Mannheim, \today \hspace{7cm} Unterschrift\\

%Abtretungs Erklaerung
\chapter*{Abtretungserkl\"arung}
\thispagestyle{empty}
Hinsichtlich meiner Studienarbeit/Bachelor-Abschlussarbeit/Diplomarbeit r\"aume ich der Universit\"at Mannheim/Lehrstuhl f{\"u}r Softwaretechnik,
Prof. Dr. Colin Atkinson, umfassende, ausschlie{\ss}liche unbefristete und
unbeschr\"ankte Nutzungsrechte an den entstandenen Arbeitsergebnissen ein.

Die Abtretung umfasst das Recht auf Nutzung der Arbeitsergebnisse in Forschung
und Lehre, das Recht der Vervielf\"altigung, Verbreitung und \"Ubersetzung sowie
das Recht zur Bearbeitung und \"Anderung inklusive Nutzung der dabei
entstehenden Ergebnisse, sowie das Recht zur Weiter\"ubertragung auf Dritte.

Solange von mir erstellte Ergebnisse in der urspr\"unglichen oder in
\"uberarbeiteter Form verwendet werden, werde ich nach Ma{\ss}gabe des Urheberrechts
als Co-Autor namentlich genannt. Eine gewerbliche Nutzung ist von dieser
Abtretung nicht mit umfasst.
\bigskip

\vspace{2.5cm}

Mannheim, \today \hspace{7cm} Unterschrift\\


\end{document}
